\documentclass[protokoll]{dvd}

\KOMAoptions{
    headwidth = 18cm,
    footwidth = 18cm,
}

% Styrelsen
\def\ordf{Daniell Cole}
\def\vice{\vak}
\def\samo{Angelika Hagberg}
\def\kassa{Owais Tabbaa}
\def\sekr{\vak}

\def\vak{\emph{vakant}}

% mötesinfo

\def\nr{3}
\def\yr{2026}
\def\ymd{2026-10-01}
\def\starttime{12:12}


\begin{document}

\title{Styrelsemöte \nr}
\subtitle{\yr}
\author{Styrelsen}
\date{\ymd}


\textbf{Datum:} \csname @date\endcsname\\
\textbf{Tid:} \starttime\\
\textbf{Plats:} Styrelserummet\\
\textbf{Styrelsemedlemmar:}
\begin{närvarande_förtroendevalda}
    \förtroendevald{Ordförande}{\ordf}{Ja}
    \förtroendevald{Vice Ordförande}{\vice}{---}
    \förtroendevald{Kassör}{\kassa}{Ja}
    \förtroendevald{SAMO}{\samo}{Ja}
    \förtroendevald{Sekreterare}{\sekr}{---}
\end{närvarande_förtroendevalda}

%\textbf{Övriga medlemmar:}
%\begin{närvarande_medlemmar}
%\medlem{Ingen}[]
%\end{närvarande_medlemmar}
\section{Öppnande av möte}

Mötet planeras öppnas av \ordf\ kl \starttime


Öppnades 12:44

\section{Runda bordet}

Runda bordet innebär att varje person berättar hur de känner sig.
Man kan till exempel berätta att man är stressad på grund av en inlämning, irriterad på sin granne, eller bara väldigt glad därför att man ligger i fas med plugget.

\newpage
\section{Formalia}

\subsection{Styrelsens beslutbarhet}

\blockquote[7 kap. 5 \S~första stycket i stadgan][]{%
    Styrelsen är endast beslutsmässig då samtliga styrelsemedlemmar har fått kallelsen till styrelsemötet och minst hälften av styrelsemedlemmarna är närvarande. Ordförande eller vice ordförande måste vara närvarande när beslut tas.
}

\blockquote[7 kap. 11 \S~första stycket i stadgan]{
Styrelsens sammanträden måste kallas minst två läsdagar innan sammanträdet.}

\blockquote[7 kap. 11 \S~sista stycket i stadgan]{
Kallelse till sammanträde får endast ske när protokollet från tidigare sammanträde är justerat och tillgängligt för medlemmarna.}

\subsubsection*{Förslag}

\begin{attsatser}
    \item Styrelsen har uppnått kraven i 7 kap. 5 § första stycket i stadgan och är därmed beslutbar.
\end{attsatser}
\subsubsection*{Beslut}
\begin{attsatser}
    \item förslaget till beslut bifalles
\end{attsatser}


\subsection{Fastställande av mötesschema}
För att styrelsen ska kunna fatta ett styrelsebeslut eller protokollföra en diskussion behöver punkten i mötesschemat där styrelsen ska fatta beslut vara inlagd eller föras in i mötesschemat senast vid den här punkten.

\subsubsection*{Förslag}

\begin{attsatser}
    \item 
    mötesschemat fastställs utan några förändringar.
    %punkt läggs till på mötesschema
\end{attsatser}
\subsubsection*{Beslut}
\begin{attsatser}
    \item förslaget till beslut bifalles
\end{attsatser}


\subsection{Val av mötessekreterare}
Antecknar det som sägs och beslutas.

\subsubsection*{Förslag}
\begin{attsatser}
    \item Owais väljs till mötessekreterare
\end{attsatser}
\subsubsection*{Beslut}
\begin{attsatser}
    \item attsats 1 bifalles
\end{attsatser}

\subsection{Val av protokolljusterare}

Protokolljusterare har till uppgift att kontrollera att protokollet i slutändan reflekterar de faktiska besluten och diskussionerna som fördes under mötet.
Utöver protokolljusteraren så ska mötesordförande och mötessekreteraren signera protokollet.
Vid styrelsemöten ska det endast vara en justerare.
Mötesordförande och mötessekreteraren kan inte vara justerare.

\subsubsection*{Förslag}
\begin{attsatser}
    \item Angelika väljs till justerare
\end{attsatser}
\subsubsection*{Beslut}
\begin{attsatser}
    \item förslaget till beslut bifalles
\end{attsatser}

\section{Rapport} 
\subsection{Ordförande}
Äntligen har banksakerna fixade och vi har nu giltig styrele enligt både banken och Göta. 

Var på FUM 1, anteckningar kommer läggas upp på DVD:Aktiva discordservern. Behövde gå tidigt från mötet pga mental hälsa. 

Färdigformulerat första utskick av de stadgaändringar som 1, göta vill att vi inför och 2, som jag tror skulle vara bra att införa, förslagen är med i protokollet på diskussionspunkter. 

\subsection{Kassör}
Äntligen tillgång till kontot.
Behövde tydligen ingen dosa, nytt från förra året. Kan nu börja jobba, men behöver nu en bokföringslektion. Insparkspengar kommer imorgon. 

\subsection{SAMO}
Inget att rapportera. 


\newpage


\section{Diskussionspunkter}
\subsection{Aspning, vad och när}
Vilka event vill vi hålla, och när vill vi ha dem?

Vi har tidigare haft pubrunda och gå till Rottan så förslagsvis göra det igen. För tillfället har Rottan stängt köket men det är på gång så förhoppningsvis är det igång i full fart någon gång snart. 

Kåris och pubrunda onsdag 4:e nov. 
Rottan fredag 20:e nov. 

\subsection{Nästa stämma}
Bara en i början av december/slutet av november eller en till emellan?

Vi tar beslut om det nästa möte. 

\subsection{Stadgaändringar}
Förslag 1, minimum 3 medlemmar i styrelsen: 
\begin{attsatser}
    \item i början av 7 kap. 3 § i stadgan införa 
    \begin{quote}
        Styrelsen kan som minst bestå av tre medlemmar varav en skall vara divisionsordförande och en divisionskassör. 
    \end{quote}
\end{attsatser}

Förslag 2, krav från Göta, styrelsemedlem behöver vara ordinarie medlem i Göta studentkår:
\begin{attsatser}
    \item ändra första stycket i 7 kap. 2 § från
    \begin{quote}
        Styrelsens medlemmar väljs av divisionsstämman. Mandatperioden är densamma som räkenskapsåret. Valbar till styrelsen är medlem i divisionen.
    \end{quote}
    till
    \begin{quote}
        Styrelsens medlemmar väljs av divisionsstämman. Mandatperioden är densamma som räkenskapsåret. Valbar till styrelsen är medlem i divisionen med ordinarie medlemskap i Göta studentkår.
    \end{quote}
\end{attsatser}

Förslag 3, att till krav från Göta, sammanfatta vilka rättigheter föreningsmedlemmar har:
\begin{attsatser}
    \item införa i 4 kap, 4 § av stadgan 
    \begin{quote}
        Medlemskap medför: 
        \begin{itemize}
            \item rätten att kalla till divisionsstämma (6 kap, 1 §),
            \item närvaro-, yttrande, förslags- och rösträtt på divisionsstämman (6 kap, 5 §),
            \item rätt att bli medlem i styrelsen (7 kap, 2 §),
            \item närvarorätt på styrelsemöten (7 kap, 6 §),
            \item samt rätten att bli medlem i kommittéer (8 kap, 4 §). 
        \end{itemize}
    \end{quote}
\end{attsatser}

Förslag 4, den stora jag främst vill ha feedback på, ny roll i styrelsen; Styrelseledamot:
\begin{attsatser}
    \item ändra 7 kap, 1§ från:
    \begin{quote}
    Styrelsen består av
		\begin{itemize}
			\item ordförande (även kallad divisionsordförande);
			\item vice ordförande;
			\item sekreterare;
			\item kassör (även kallad divisionsskassör);
			\item SAMO; och
            \end{itemize}
    \end{quote}
    till:
    \begin{quote}
        Styrelsens ordinarie medlemmar består av 
		\begin{itemize}
			\item ordförande (även kallad divisionsordförande);
			\item vice ordförande;
			\item sekreterare;
			\item kassör (även kallad divisionsskassör);
			\item SAMO; och
            \end{itemize}
            ytterligare styrelsemedlemmar består av
            \begin{itemize}
                \item ledamöter (även kallat    styrelseledamöter).
            \end{itemize}
    \end{quote}
    
    \item införa i 7 kap. 2 §:
    \begin{quote}
        Ordinarie styrelsemedlem kan inte samtidigt vara styrelseledamot.
    \end{quote}
    
    \item uppdatera 7 kap. 3 § från:
    \begin{quote}
        Styrelsen kan högst bestå av fem medlemmar.
    \end{quote}
    till 
    \begin{quote}
        Styrelsen kan högst bestå av fem ordinarie medlemmar. 
    \end{quote}
    
    \item 
    ändra första stycket i 7 kap. 5 § från:
    \begin{quote}
        Styrelsen är endast beslutsmässig då samtliga styrelsemedlemmar har fått kallelsen till styrelsemötet och minst hälften av styrelsemedlemmarna är närvarande.
        Ordförande eller vice ordförande måste vara närvarande när beslut tas.
    \end{quote}
    till:
    \begin{quote}
        Styrelsen är endast beslutsmässig då samtliga styrelsemedlemmar har fått kallelsen till styrelsemötet och minst hälften av styrelsens ordinarie medlemmar är närvarande.
        Ordförande eller vice ordförande måste vara närvarande när beslut tas.
    \end{quote}
\end{attsatser}

\newpage
\section{Besultspunkter}
\subsection{Roller i styrelsen}
Då styrelsen konstituerar sig själv måste styrelsen nu välja roller.
\subsection*{Förslag till beslut:}
\begin{attsatser}
    \item Angelika väljs in som SAMO för verksamhetsåret 2026.
\end{attsatser}
\subsection*{Beslut}
\begin{attsatser}
    \item attsats 1 bifalles.
\end{attsatser}


\subsection{Fylla första hjälpen lådan}
Första hjälpen lådan är tom, den behöver fyllas. Kostar 527,50 kr inklusive moms. Finns att köpa här.
Men undersök först om Akademiska hus kan stå för det. 

\url{https://www.forstahjalpencentrum.se/produkter/refill-till-cederroth-forsta-hjalpen-tavla/?gad_source=1&gad_campaignid=23240897969&gclid=Cj0KCQjwteTUBhD4ARIsAEYjs3pM1Oe38jaU_7z9j_ijrq3bziz2LZLCrdh4via3tbx8EJ_DD_661XEaArsYEALw_wcB}

\subsubsection*{Förslag till beslut}
\begin{attsatser}
    \item Owais undersöker om Akademiska hus kan stå för påfyllnaden, men om de inte gör det införskaffa första hjälpen kit för 527,50 kr. 
\end{attsatser}
\subsubsection*{Beslut}
\begin{attsatser}
    \item Attsats 1 bifalles.
\end{attsatser}

\newpage
\section{Per Capsulam-beslut}
\blockquote[]{
När per capsulam-beslut har fattats måste beslutet dokumenteras,\\
\ \cdot\ \ på divisionens Discordserver som är tillgänglig alla studenter;\\
\ \cdot\ \ i protokollet för det första styrelsesammanträdet efter beslut har fattats.\\

Dokumentationen ska innehålla\\
\ \cdot\ \ vilka att-satser som har bifallits;\\
\ \cdot\ \ omröstningsperioden.\\

Om dessa krav inte uppfylls upphävs beslutet.
}
\textit{Inga per Capsulam-beslut har tagits}
%\subsubsection*{Bifallna attsatser}
%\subsubsection*{Omröstningsperioden}

%\newpage
\section{Avslutande av möte}

\subsection{Nästa möte}
2026-10-15, 11:11

\subsection{Mötets avslutande}
Mötet avslutades kl 13:44 

\styrelsesignaturer

\end{document}